%%%%%%%%%%%%%%%%%%%%%%%%%%%%%%%%%%%%%%%%%%%%%%%%%%%%%%%%%%%%%%%%%%%%%%%%%%%
% Lehreinheit 4: Quellenkritik
% Quelle: WuZ-Buch, Kap. 15 (LE 4), Kap. 4 (Confirmation Bias),
%         Kap. 5 (Alief, Vanishing Ball), Kap. 15b (Entlarvung:
%         Prebunking, laterales Lesen, Forschungsbetrug)
% Staunmoment: Zufallspunktwolken (Apophaenie) und "schoene Faelschung
%         vs. haessliches Original"; optional "Das letzte Abendmahl"
%         (Begleitbuch, Anhang K)
% Kernbegriff: Apophaenie / Quellenkritik
%%%%%%%%%%%%%%%%%%%%%%%%%%%%%%%%%%%%%%%%%%%%%%%%%%%%%%%%%%%%%%%%%%%%%%%%%%%

% ---- Punktwolken fuer den Staunmoment (reproduzierbar per Seed) -----------
% \LEvierZufall{Seed}{Beschriftung}: 120 gleichverteilte Zufallspunkte
\newcommand{\LEvierZufall}[2]{%
\begin{tikzpicture}[x=1cm,y=1cm]
  \draw[rwugray!70] (0,0) rectangle (4,4);
  \pgfmathsetseed{#1}%
  \foreach \i in {1,...,120}{%
    \pgfmathsetmacro{\px}{0.1+3.8*rnd}%
    \pgfmathsetmacro{\py}{0.1+3.8*rnd}%
    \fill[rwuviolet] (\px,\py) circle (1.4pt);}
  \node[below,font=\sffamily\bfseries] at (2,-0.05) {#2};
\end{tikzpicture}}
% \LEvierGitter{Seed}{Beschriftung}: 11x11-Gitter mit leichtem Zittern
\newcommand{\LEvierGitter}[2]{%
\begin{tikzpicture}[x=1cm,y=1cm]
  \draw[rwugray!70] (0,0) rectangle (4,4);
  \pgfmathsetseed{#1}%
  \foreach \i in {0,...,10}{\foreach \j in {0,...,10}{%
    \pgfmathsetmacro{\px}{0.2+0.36*\i+0.26*(rnd-0.5)}%
    \pgfmathsetmacro{\py}{0.2+0.36*\j+0.26*(rnd-0.5)}%
    \fill[rwuviolet] (\px,\py) circle (1.4pt);}}
  \node[below,font=\sffamily\bfseries] at (2,-0.05) {#2};
\end{tikzpicture}}
% \LEvierLinie{Seed}{Beschriftung}: 95 Zufallspunkte + 25 Punkte nahe
% der Geraden y = 0.6x + 0.9 (schwaches echtes Muster)
\newcommand{\LEvierLinie}[2]{%
\begin{tikzpicture}[x=1cm,y=1cm]
  \draw[rwugray!70] (0,0) rectangle (4,4);
  \pgfmathsetseed{#1}%
  \foreach \i in {1,...,95}{%
    \pgfmathsetmacro{\px}{0.1+3.8*rnd}%
    \pgfmathsetmacro{\py}{0.1+3.8*rnd}%
    \fill[rwuviolet] (\px,\py) circle (1.4pt);}
  \foreach \i in {1,...,25}{%
    \pgfmathsetmacro{\px}{0.2+3.6*rnd}%
    \pgfmathsetmacro{\py}{0.6*\px+0.9+0.24*(rnd-0.5)}%
    \fill[rwuviolet] (\px,\py) circle (1.4pt);}
  \node[below,font=\sffamily\bfseries] at (2,-0.05) {#2};
\end{tikzpicture}}

\nsl{In der dritten Einheit haben die Studierenden Literatur gefunden. Jetzt
müssen sie entscheiden, welcher sie trauen. Die Einheit beginnt mit einer
Erfahrung, die unangenehm ist und gerade deshalb hängen bleibt: Das eigene
Gehirn findet in reinem Zufall Muster, und es hält einen gut gestalteten
Text für glaubwürdiger als einen schlecht gesetzten, ohne den Inhalt
geprüft zu haben. Quellenkritik ist die Methode, die gegen beides hilft.
Sie wird hier als angewandte Skepsis gegen die eigene Mustererkennung
eingeführt: Kriterien und Qualitätsstufen wissenschaftlicher Literatur,
die Grenzen von Plausibilitätsurteilen und Zitationsmetriken, und als
Werkzeug für den Alltag das laterale Lesen nach dem SIFT-Schema.
\newpage}
\os{\begin{description}
\item[{\rot\bf Heute}]~\\[-\lblskip]
\bi
  \ite Muster im Rauschen: Apophänie
  \ite schöne Fälschung, hässliches Original
  \ite fünf Kriterien und vier Qualitätsstufen
  \ite Impact Factor und h-Index: was sie nicht sagen
  \ite laterales Lesen mit SIFT, Prebunking
\ei
\end{description}
\newpage}

\begin{vorbereitung}[Punktwolken und zwei Artikel]
\textbf{Eine Woche vorher}
\begin{checkliste}
  \item Die Punktwolken liegen als Folien vor (sie werden beim Übersetzen
    aus festen Zufallszahlen erzeugt und sehen jedes Mal gleich aus). Die
    Folien einmal durchklicken und die Lösung (Lehrhinweis nach dem
    Staunmoment) nicht auf den Beamer legen.
  \item \textbf{Die \enquote{schöne Fälschung} selbst bauen.} Ein fiktiver
    Nachrichtenartikel zu einem \emph{harmlosen} Thema (Beispiel: \enquote{Wer
    morgens zehn Minuten jongliert, merkt sich Vokabeln besser}). Er
    soll alle Merkmale tragen, die Glaubwürdigkeit \emph{erzeugen}, ohne sie
    zu belegen: saubere Typographie, Foto, Diagramm, ein Experte mit Titel,
    eine präzise Zahl (\enquote{um 23,7\,\%}), der Satz \enquote{wie eine neue
    Studie zeigt} ohne nachprüfbaren Verweis. Keine Logos, Namen oder
    Domains realer Medien und keine realen Personen verwenden. In der
    Fußzeile klein, aber lesbar: \enquote{Fiktives Lehrmaterial, LE~4}.
  \item \textbf{Das \enquote{hässliche Original} auswählen.} Ein echter,
    begutachteter Fachartikel mit unattraktivem Layout: eingescannte
    Schreibmaschinenseiten, zweispaltiger Satz ohne Abbildungen, lange
    Tabellen. Kriterien: frei zugänglich, Befund seither nicht
    zurückgezogen, Thema für Erstsemester verständlich, Abstract auf einer
    Folie lesbar. Gut geeignet sind ältere psychologische Klassiker, die im
    Begleitbuch ausgewertet werden, etwa \textcite{forer1949fallacy} oder
    \textcite{wason1960failure}; die Rechte für die Projektion im Hörsaal
    vorher klären.
  \item Beide Artikel als Bildschirmfoto auf je eine Folie legen, gleich
    groß, ohne Quellenangabe. Die Quellenangabe des echten Artikels liegt
    auf einer dritten Folie für die Auflösung.
  \item \textbf{Drei Quellen für die SIFT-Übung festlegen} (Pflicht, sie
    stehen nicht im Skript): eine Peer-Review-Studie, ein Zeitungsartikel,
    der über eine Studie berichtet, und ein Wikipedia-Eintrag. Am
    ergiebigsten ist eine \emph{Kette}: Zeitungsartikel und
    Wikipedia-Eintrag berufen sich auf dieselbe Studie, sodass die Gruppe
    in Schritt~T die Behauptung bis zum Original verfolgen kann. Bei der
    Auswahl darauf achten, dass der Zeitungsartikel die Studie zuspitzt
    (eine Zahl, die im Original vorsichtiger formuliert ist), und den
    Wikipedia-Eintrag als Versionslink (\enquote{Permanenter Link})
    ausgeben, damit alle Gruppen dieselbe Fassung prüfen.
\end{checkliste}

\textbf{Handout und Übungsblatt}
\begin{checkliste}
  \item Handout \enquote{Quellenkritik} und Übungsblatt LE~4 je einmal pro
    Person kopieren (Druckvorlage \texttt{handouts.pdf}, Seiten zu LE~4).
  \item Die drei SIFT-Quellen als Links (QR-Code oder Kurz-URL) auf eine
    Folie und auf das Lernportal legen; zusätzlich je Kleingruppe einen
    Ausdruck der drei Quellen (Bibliothekszugang für die Studie prüfen).
  \item Für die Sicherung ein A3-Blatt oder eine Tafelhälfte mit der
    Überschrift \enquote{Quellenkritik in fünf Schritten}.
\end{checkliste}

\textbf{Am Tag selbst}
\begin{checkliste}
  \item WLAN prüfen: Die SIFT-Übung braucht in jeder Gruppe mindestens ein
    Gerät mit Internetzugang.
  \item Tafel teilen: links \enquote{Muster, die wir gesehen haben}, rechts
    \enquote{Woran wir Glaubwürdigkeit festgemacht haben}.
  \item Zwei Studierende vorab bitten, im Blitzlicht je eine annotierte
    Quelle aus LE~3 vorzustellen.
\end{checkliste}

\textbf{Optionale Erweiterung: \enquote{Das letzte Abendmahl}.} Wer das
Kunststück aus dem Begleitbuch beherrscht, kann den Staunmoment damit
eröffnen (\buch{Anhang~K}; Kunststück nach \textcite{payne_davinci}). Ein
gefaltetes Blatt zeigt nacheinander mehrere Bilder, zuletzt Leonardos
Abendmahl und die Mona Lisa mit der gewählten Karte. Die augenzwinkernde
\enquote{Symbolanalyse} des Gemäldes ist eine Miniatur über Apophänie,
motiviertes Sehen und Verschwörungsdenken und leitet direkt zu den
Punktwolken über. Methode und Einstudierung nach der Vorlage im
Begleitbuch; hier bewusst nicht beschrieben.

\textbf{Risiken und Plan B.} Manche Studierende ahnen bei den Punktwolken
die Pointe und sagen \enquote{das ist Zufall}. Das ist willkommen: Dann
die Frage verschärfen (\enquote{Welche der drei ist \emph{am wenigsten}
zufällig?}) und die zweite Folie zeigen; dort irren erfahrungsgemäß auch
die Vorsichtigen. Wenn beim Artikelvergleich die Mehrheit sofort das
Original wählt, nachfragen, \emph{woran} sie es erkannt hat. Meist sind es
wieder Oberflächenmerkmale (\enquote{wirkt zu glatt}), also ebenfalls eine
Heuristik und keine Prüfung. Die fiktive Fälschung nach der Stunde nicht
digital weitergeben.
\end{vorbereitung}

\begin{ablauf}
\begin{description}
\item[0--5 · Blitzlicht.] Zwei annotierte Quellen aus LE~3 werden
  vorgestellt: \enquote{Woran habe ich Qualität festgemacht?} Antworten an
  die rechte Tafelhälfte.
\item[5--20 · Staunmoment.] Punktwolken: \enquote{Welche Muster sehen
  Sie?}, dann \enquote{Welche ist Zufall?} Danach der Vergleich
  \enquote{schöne Fälschung} und \enquote{hässliches Original} mit
  Handabstimmung. Auflösung beider Teile; Staunen über die eigene
  Fehlleistung.
\item[20--50 · Theorie-Input.] Kriterien und Qualitätsstufen der
  Quellenkritik; Apophänie und Bestätigungsfehler; Plausibilitätsheuristiken
  und ihre Grenzen; Metriken; laterale Recherche und Prebunking.
\item[50--80 · Anwendung (SIFT).] Kleingruppen prüfen die drei
  vorgegebenen Quellen nach dem SIFT-Schema (Übungsblatt, Aufgabe~2).
\item[80--90 · Sicherung.] Gemeinsame Checkliste \enquote{Quellenkritik in
  fünf Schritten}; Diskussionsauftrag zur Zitierhäufigkeit als
  Qualitätsindikator.
\end{description}
\end{ablauf}

%==============================================================================
\section{Staunmoment: Muster im Rauschen}
%==============================================================================
\begin{staunbox}{Welche Muster sehen Sie?}
Drei Punktwolken, je 120 Punkte. Sehen Sie sich jede zehn Sekunden lang an.

\emph{Notieren Sie alles, was Ihnen auffällt: Häufungen, Linien, Lücken,
Figuren.}
\end{staunbox}
\os{\newpage}

\begin{description}
\item[{\rot\bf Welche Muster sehen Sie?}]~\\[-\lblskip]
\begin{center}
\resizebox{0.95\textwidth}{!}{%
\LEvierZufall{4711}{A}\hspace{6mm}%
\LEvierZufall{1848}{B}\hspace{6mm}%
\LEvierZufall{2604}{C}}
\end{center}
\os{\newpage}

\item[{\rot\bf Und jetzt: Welche ist Zufall?}]~\\[-\lblskip]
\begin{center}
\resizebox{0.95\textwidth}{!}{%
\LEvierGitter{77}{D}\hspace{6mm}%
\LEvierZufall{9001}{E}\hspace{6mm}%
\LEvierLinie{314}{F}}
\end{center}
\bi
  \ite eine Wolke ist reiner Zufall, eine ist geordnet, eine enthält ein
    echtes Muster
\ei
\os{\newpage}
\end{description}

\begin{lehrhinweis}[Auflösung der Punktwolken]
A, B und C sind gleichverteilte Zufallspunkte (feste Startwerte, damit
die Bilder reproduzierbar sind). Jede Häufung, Lücke oder \enquote{Linie},
die die Gruppe findet, ist Rauschen. Auf der zweiten Folie ist E reiner
Zufall; D ist ein regelmäßiges Gitter mit leichtem Zittern, F enthält 25
Punkte nahe einer schrägen Geraden. Häufig wird D für den Zufall gehalten,
weil die Punkte so \enquote{gleichmäßig verstreut} wirken, und E für
geordnet, weil echter Zufall Klumpen bildet. Das Ergebnis vor der
Auflösung per Handzeichen festhalten. Die Pointe: Unser Gefühl für
\enquote{zufällig} ist falsch kalibriert, und ein schwaches echtes Muster
(F) ist mit bloßem Auge kaum von den Scheinmustern in A--C zu
unterscheiden. Ob ein Muster echt ist, entscheidet nicht der Blick, sondern
ein Verfahren (Statistik, Replikation, LE~8).
\end{lehrhinweis}

\begin{description}
\item[{\rot\bf Zwei Artikel}]~\\[-\lblskip]
\bi
  \ite Sie sehen zwei Ausschnitte, beide 30 Sekunden lang
  \ite Welcher ist glaubwürdiger? Handzeichen
  \ite Woran haben Sie das festgemacht? Drei Merkmale notieren
\ei
\nsl{Die beiden Folien mit den Artikeln folgen hier (siehe Vorbereitung).
Erst nach der Abstimmung wird aufgelöst: Der elegante Artikel ist frei
erfunden und von der Lehrperson für diese Stunde gebaut, der unscheinbare
ist eine begutachtete Originalarbeit. Die notierten Merkmale gehen an die
rechte Tafelhälfte und werden im Theorieteil geordnet: Fast alle betreffen
die Form (Layout, Bild, Zahl, Titel), kaum eines die Herkunft oder die
Methode.}
\os{\newpage}

\item[{\rot\bf Was ist gerade passiert?}]~\\[-\lblskip]
\bi
  \ite in reinem Rauschen haben wir Muster gesehen
  \ite Ordnung hielten wir für Zufall
  \ite einen erfundenen Text hielten viele für glaubwürdiger als eine
    echte Studie
  \ite unsere Kriterien betrafen die \emph{Form}, nicht die \emph{Herkunft}
\ei
\os{\newpage}
\end{description}

%==============================================================================
\section{Apophänie und Bestätigungsfehler}
%==============================================================================
\nsl{Die Neigung, in zufälligen Daten bedeutungsvolle Zusammenhänge zu
erleben, hat der Psychiater Klaus Conrad \emph{Apophänie} genannt. Er
beschrieb damit ursprünglich eine Phase beginnender Psychosen, in der den
Betroffenen alles bedeutsam und auf sie bezogen erscheint
\parencite{conrad1958schizophrenie}. Heute bezeichnet der Begriff eine
allgemeine Eigenschaft der Wahrnehmung. Sie verstärkt sich unter
Kontrollverlust: In den Experimenten von \textcite{whitson2008lacking}
sahen Versuchspersonen, deren Kontrollgefühl man zuvor geschwächt hatte,
häufiger Bilder in Rauschmustern und häufiger Verschwörungen.

Das passt zu dem, was das Begleitbuch über Wahrnehmung als
Hypothesenbildung sagt (\buch{Kap.~5}). Bei der \emph{Vanishing Ball
Illusion} gaben zwei Drittel der Versuchspersonen an, einen Ball
hochfliegen und verschwinden gesehen zu haben, der die Hand nie verlassen
hatte \parencite{kuhn2016}. Das Gehirn ergänzt, was es erwartet. In
Rauschdaten ergänzt es Muster, in Texten ergänzt es Glaubwürdigkeit.

Der \emph{Bestätigungsfehler} setzt eine Stufe später an: Er filtert die
Bewertung der Evidenz, die wir bereits gesehen haben
\parencite{nickerson1998ubiquitous}. Wason zeigte das mit einer
Zahlenfolge, deren Regel die Versuchspersonen erraten sollten; sie testeten
fast nur Beispiele, die ihre Vermutung bestätigen konnten
\parencite{wason1960failure}. Für die Literaturarbeit heißt das: Wir
finden Quellen besonders überzeugend, die bestätigen, was wir schon
glauben, und prüfen die anderen strenger.}
\begin{description}
\item[{\rot\bf Apophänie}]~\\
\defn{apophaenie}{\textbf{Apophänie} ist die Neigung, in zufälligen oder
bedeutungslosen Daten Muster und Zusammenhänge zu erleben
\parencite{conrad1958schizophrenie}.}
\bi
  \ite die Wahrnehmung bildet Hypothesen, ergänzt Fehlendes
    (\buch{Kap.~5})
  \ite Vanishing Ball: zwei Drittel \enquote{sehen} einen Ball, der nie
    geworfen wurde \parencite{kuhn2016}
  \ite unter Kontrollverlust sehen Menschen mehr Scheinmuster
    \parencite{whitson2008lacking}
\ei
\os{\newpage}

\item[{\rot\bf Der Bestätigungsfehler}]~\\[-\lblskip]
\bi
  \ite wir prüfen Belege, die uns recht geben, weniger streng
    \parencite{nickerson1998ubiquitous}
  \ite Probe: Die Folge 8 -- 10 -- 12 folgt einer Regel. Welche?
  \ite fast alle testen nur Beispiele, die ihre Vermutung
    \emph{bestätigen} könnten \parencite{wason1960failure}
  \ite für die Recherche: Wer nach Bestätigung sucht, findet sie
\ei
\nsl{Die Probe mit 8 -- 10 -- 12 ist im Begleitbuch als Hörsaalexperiment
beschrieben (\buch{Kap.~4}): Die meisten vermuten \enquote{gerade Zahlen}
oder \enquote{immer plus zwei} und nennen 14. Die Regel lautet
\enquote{aufsteigende Zahlen}. Sie hätten sie nur gefunden, wenn sie eine
Folge getestet hätten, die ihrer Vermutung widerspricht. Wer in LE~3 eine
Suchanfrage formuliert hat, die nur die erwartete Antwort finden kann, hat
denselben Fehler gemacht.}
\os{\newpage}

\item[{\rot\bf Quellenkritik als Gegenmittel}]~\\[-\lblskip]
\bi
  \ite gegen Apophänie: Muster erst glauben, wenn ein \emph{Verfahren} sie
    bestätigt
  \ite gegen den Bestätigungsfehler: dieselben Kriterien an
    \emph{alle} Quellen legen, auch an die bequemen
  \ite gezielt nach der stärksten Gegenposition suchen
\ei
\os{\newpage}
\end{description}

\begin{merkbox}
Das Gefühl, ein Muster zu sehen, ist kein Beleg dafür, dass es da ist. Das
Gefühl, einer Quelle trauen zu können, ist kein Beleg für ihre Qualität.
Beides wird geprüft, nicht empfunden.
\end{merkbox}
\os{\newpage}

%==============================================================================
\section{Kriterien der Quellenkritik}
%==============================================================================
\nsl{Die Geschichtswissenschaft hat die Quellenkritik als Methode
entwickelt; die empirischen Wissenschaften haben sie um die Frage nach der
Methode erweitert. Für wissenschaftliches Arbeiten genügen fünf
Kriterien, die sich in jeder Disziplin anwenden lassen. Keines entscheidet
allein, und keines ist mit einem Blick auf das Layout zu beantworten.}
\begin{description}
\item[{\rot\bf Was ist Quellenkritik?}]~\\
\defn{quellenkritik}{\textbf{Quellenkritik} ist die methodische Prüfung,
wer eine Information mit welcher Absicht, auf welchem Weg und mit welcher
Prüfung erzeugt hat, um zu entscheiden, wie viel Gewicht sie in der
eigenen Argumentation tragen darf.}
\bi
  \ite das Ergebnis ist ein \emph{Gewicht}, kein Ja oder Nein
  \ite geprüft wird die Quelle, nicht ihr Aussehen
\ei
\os{\newpage}

\item[{\rot\bf Fünf Kriterien}]~\\[-\lblskip]
\begin{center}
\small
\begin{tabular}{@{}p{0.2\textwidth}p{0.37\textwidth}p{0.33\textwidth}@{}}
\toprule
\textbf{Kriterium} & \textbf{Prüffrage} & \textbf{Warnsignal} \\
\midrule
Herkunft & Wer spricht? Welche Expertise, welche Institution? & kein Autor, kein Impressum \\
Intention & Wozu wurde der Text geschrieben? Wer bezahlt? & verkauft etwas, Interessenkonflikt verschwiegen \\
Methode & Wie wurde das Ergebnis gewonnen? & \enquote{Studien zeigen} ohne Angabe welcher \\
Peer Review & Haben unabhängige Fachleute geprüft? & Selbstverlag, Raubverlag \\
Aktualität & Ist das noch Stand der Forschung? & widerrufen, überholt, nie repliziert \\
\bottomrule
\end{tabular}
\end{center}
\os{\newpage}
\end{description}

\nsl{\textbf{Herkunft.} Gefragt wird nach der Person und der Institution,
nach deren Fachkompetenz für \emph{genau diese} Frage und nach der
Nachprüfbarkeit dieser Angaben. Ein Nobelpreis in Chemie macht nicht zum
Fachmann für Ernährung. \textbf{Intention.} Informieren, überzeugen,
verkaufen und unterhalten sind legitime Absichten, aber sie verändern das
Gewicht einer Aussage. Eine Pressemitteilung eines Herstellers über das
eigene Produkt ist eine Quelle, aber eine parteiische. \textbf{Methode.}
Hier liegt der Kern der wissenschaftlichen Quellenkritik: Wie viele
Personen wurden untersucht, gab es eine Kontrollgruppe, sind die Daten
offengelegt? Die Frage wird in LE~6 und LE~8 ausgebaut. \textbf{Peer
Review.} Die Begutachtung vor der Veröffentlichung ist eine wichtige,
aber begrenzte Prüfung (siehe unten). \textbf{Aktualität.} Ein alter Text
ist nicht schlechter als ein neuer; aber man muss wissen, ob sein Befund
inzwischen widerrufen, bestätigt oder überholt ist.}

\begin{description}
\item[{\rot\bf Was Peer Review leistet -- und was nicht}]~\\[-\lblskip]
\bi
  \ite Gutachter prüfen Logik, Methode, Einordnung, Darstellung
  \ite sie sehen in der Regel \emph{nicht} die Rohdaten
  \ite Stapel: mindestens 55 Publikationen mit erfundenen Daten, entdeckt
    von Mitarbeitern, nicht von Gutachtern \parencite{levelt2012flawed}
  \ite Bik: problematische Bildduplikate in 3,8\,\% von 20.621 Aufsätzen
    \parencite{bik2016prevalence}
  \ite Peer Review ist eine Mindestprüfung, kein Wahrheitssiegel
\ei
\nsl{Beide Fälle sind im Begleitbuch ausführlich dargestellt
(\buch{Kap.~15b}). Bemerkenswert ist der zweite Teil der Geschichte: Die
Fehler wurden \emph{innerhalb} der Wissenschaft aufgedeckt, durch
nachträgliche Prüfung. Das ist kein Argument gegen Peer Review, sondern
eines dafür, Begutachtung als laufenden Prozess zu verstehen, der mit der
Veröffentlichung nicht endet. Für die Quellenkritik folgt daraus, dass
\emph{nach} der Veröffentlichung geprüft werden muss: Gibt es
Korrekturen, Kommentare, einen Widerruf?}
\os{\newpage}
\end{description}

%==============================================================================
\section{Qualitätsstufen wissenschaftlicher Literatur}
%==============================================================================
\nsl{Neben den Kriterien für die einzelne Quelle braucht man eine Vorstellung
davon, welche \emph{Art} von Text man vor sich hat. Eine
\emph{Primärstudie} berichtet eigene Daten. Ein \emph{Review} fasst den
Stand mehrerer Studien zusammen, entweder erzählend (narrativ) oder nach
einem vorab festgelegten, offengelegten Suchprotokoll (systematisch). Eine
\emph{Meta-Analyse} fasst die Effektgrößen vieler Studien statistisch
zusammen. \emph{Graue Literatur} umfasst alles, was außerhalb der
Verlagswege erscheint: Preprints, Berichte von Behörden und Verbänden,
Abschlussarbeiten, Tagungsbeiträge. Zeitungsartikel und
Lexikoneinträge sind Tertiärquellen: Sie berichten über Forschung, sie
sind keine.

Für Fragen nach der Wirksamkeit einer Maßnahme wird diese Rangfolge
häufig als Pyramide dargestellt. Die Pyramide ist eine Faustregel. Eine
Meta-Analyse aus schwachen Einzelstudien bleibt schwach, und eine
qualitative Studie ist für eine Frage nach Erleben und Bedeutung die
angemessene Primärquelle, auch wenn sie in der Pyramide weit unten
stünde (LE~9). \textcite{ioannidis2005false} hat gezeigt, warum sich
gerade in Feldern mit kleinen Stichproben, kleinen Effekten und vielen
getesteten Hypothesen viele publizierte Einzelbefunde nicht halten. Daraus
folgt: Einer einzelnen Studie gibt man weniger Gewicht als einem
Befund, der unabhängig repliziert wurde.}
\begin{description}
\item[{\rot\bf Vier Qualitätsstufen}]~\\[-\lblskip]
\begin{center}
\resizebox{0.82\textwidth}{!}{%
\begin{tikzpicture}[font=\sffamily]
  \fill[rwuviolet!70] (-2.2,4.5) -- (2.2,4.5) -- (0,6.6) -- cycle;
  \fill[rwuviolet!45] (-3.8,3) -- (3.8,3) -- (2.2,4.5) -- (-2.2,4.5) -- cycle;
  \fill[rwucyan40!70] (-5.4,1.5) -- (5.4,1.5) -- (3.8,3) -- (-3.8,3) -- cycle;
  \fill[rwugray!25] (-7,0) -- (7,0) -- (5.4,1.5) -- (-5.4,1.5) -- cycle;
  \node[white,font=\sffamily\small] at (0,4.9) {\bfseries Meta-Analyse};
  \node at (0,3.75) {\bfseries systematisches Review};
  \node at (0,2.25) {\bfseries Primärstudie (begutachtet)};
  \node[align=center] at (0,0.75) {\bfseries graue Literatur\\[-0.5mm]\small Preprint, Bericht, Abschlussarbeit};
  \draw[->,thick,rwugray] (7.6,0.2) -- node[right,align=left,font=\sffamily\small]{mehr Studien\\zusammengefasst} (7.6,6.2);
  \node[rwugray,font=\sffamily\small] at (0,-0.45) {darunter: Zeitung, Lexikon, Blog (berichten über Forschung)};
\end{tikzpicture}}
\end{center}
\os{\newpage}

\item[{\rot\bf Grenzen der Pyramide}]~\\[-\lblskip]
\bi
  \ite gilt vor allem für Fragen nach Wirksamkeit
  \ite eine Meta-Analyse schwacher Studien bleibt schwach
  \ite für Fragen nach Erleben ist die qualitative Studie die richtige Primärquelle (LE~9)
  \ite viele Einzelbefunde halten nicht \parencite{ioannidis2005false}
  \ite Faustregel: repliziert vor einzeln, Original vor Bericht
\ei
\os{\newpage}
\end{description}

%==============================================================================
\section{Plausibilität und ihre Grenzen}
%==============================================================================
\nsl{Im Alltag entscheiden wir über Glaubwürdigkeit fast immer mit
Heuristiken: Klingt das plausibel? Sieht das seriös aus? Habe ich das schon
öfter gehört? Solche Faustregeln sind nicht dumm. Sie funktionieren dort,
wo die Umwelt ehrlich ist, wo also Professionalität, Wiederholung und
Plausibilität tatsächlich mit Wahrheit zusammenhängen
\parencite{gigerenzertodd1999}. Sie versagen dort, wo jemand die Hinweisreize
kennt und gezielt erzeugt. Genau das tut der Zauberkünstler auf der Bühne,
deklariert und zur Unterhaltung, und genau das tut die Desinformation,
undeklariert (\buch{Kap.~15b}).

Drei Befunde zeigen, wie wenig die Hinweisreize mit dem Inhalt zu tun
haben müssen. Aussagen werden häufiger für wahr gehalten, wenn sie leicht
lesbar dargeboten werden \parencite{reber1999fluency}. Sie werden für
wahrer gehalten, wenn man sie schon einmal gehört hat, unabhängig davon,
ob sie stimmen \parencite{hasher1977frequency}. Und ein mitreißender
Vortrag ohne sinnvollen Inhalt kann von einem Fachpublikum glänzend
bewertet werden, der sogenannte Dr.-Fox-Effekt
\parencite{naftulin1973drfox,peerbabad2014drfox}. Die schöne Fälschung im
Staunmoment hat alle drei Hebel benutzt.}
\begin{description}
\item[{\rot\bf Woran wir Glaubwürdigkeit festmachen}]~\\[-\lblskip]
\bi
  \ite \textbf{Lesbarkeit}: leicht Lesbares wirkt wahrer
    \parencite{reber1999fluency}
  \ite \textbf{Wiederholung}: schon Gehörtes wirkt wahrer
    \parencite{hasher1977frequency}
  \ite \textbf{Auftreten}: ein guter Vortrag wirkt kompetent, auch ohne
    Inhalt \parencite{naftulin1973drfox}
  \ite \textbf{Präzision}: eine genaue Zahl wirkt gemessen
\ei
\os{\newpage}

\item[{\rot\bf Die präzise Zahl}]~\\[-\lblskip]
\bi
  \ite \enquote{Die Tischdecke in Leonardos \emph{Abendmahl} enthält 52
    Rechtecke.}
  \ite klingt gezählt, ist aber erfunden
  \ite belegbar sind andere Angaben: 13 Figuren, Jesus in der Mitte,
    drei Fenster, entstanden 1495--1498
  \ite Quellenkritik fragt: \textbf{Wer hat das gezählt? Wo steht das?}
\ei
\nsl{Das Beispiel stammt aus dem Umfeld des Kunststücks \enquote{Das
letzte Abendmahl} (\buch{Anhang~K}). Die Zahl 52 passt verführerisch zu
einem Kartenspiel und wird deshalb gern weitererzählt. Präzision ist ein
Oberflächenmerkmal wie gutes Layout: Sie kostet den Erzeuger nichts und
signalisiert dem Leser Messung. Prüfen lässt sich eine solche Zahl nur, indem
man ihre Spur bis zu jemandem zurückverfolgt, der sie erhoben hat. Endet
die Spur bei einer Behauptung ohne Quelle, trägt die Zahl kein Gewicht,
gleich wie oft sie wiederholt wird.}
\os{\newpage}

\item[{\rot\bf Wann Heuristiken taugen}]~\\[-\lblskip]
\bi
  \ite Faustregeln sind schnell und oft gut \parencite{gigerenzertodd1999}
  \ite sie versagen, wo jemand die Hinweisreize gezielt erzeugt
  \ite Zauberkunst und Desinformation nutzen dieselben Hebel, nur der
    Vertrag mit dem Publikum ist verschieden (\buch{Kap.~15b})
  \ite Konsequenz: Plausibilität ist ein Anlass zu prüfen, kein Ersatz
    dafür
\ei
\os{\newpage}
\end{description}

%==============================================================================
\section{Impact Factor, h-Index und was sie nicht aussagen}
%==============================================================================
\nsl{In der Wissenschaft selbst gibt es ein Gegenstück zur Plausibilitätsheuristik:
die Zitationsmetrik. Der \emph{Impact Factor} einer Zeitschrift ist die
Zahl der Zitate, die ihre Artikel aus den beiden Vorjahren in einem Jahr
erhalten, geteilt durch die Zahl dieser Artikel. Eugene Garfield hat ihn
als Hilfsmittel entwickelt, um Zeitschriften für den Science Citation Index
auszuwählen, und selbst davor gewarnt, damit einzelne Forschende zu
bewerten \parencite{garfield2006history}. Der \emph{h-Index} einer Person
ist die größte Zahl $h$, für die sie $h$ Arbeiten mit je mindestens $h$
Zitaten hat \parencite{hirsch2005index}.

Beide Zahlen messen Aufmerksamkeit, nicht Richtigkeit.
\textcite{seglen1997impact} hat gezeigt, dass die Zitate innerhalb einer
Zeitschrift stark ungleich verteilt sind: Wenige Artikel bringen den
Großteil, die meisten liegen weit unter dem Mittelwert. Der Impact Factor
sagt deshalb wenig über den einzelnen Artikel. Zitate folgen außerdem der
Reputation: Bekannte Namen werden für vergleichbare Arbeit häufiger
zitiert, ein Effekt, den Merton nach dem Matthäus-Evangelium benannt hat
\parencite{merton1968matthew}. Und zitiert wird auch, was man widerlegt.
Stapels \emph{Science}-Artikel von 2011 \parencite{stapel2011chaos} ist
seit dem Auffliegen der Fälschung vor allem als Beispiel für Betrug von
Interesse; jede solche Erwähnung zählt als Zitat. Die \emph{San Francisco
Declaration on Research Assessment} \parencite{dora2012} und das
\emph{Leiden Manifesto} \parencite{hicks2015leiden} fordern deshalb,
Forschung nach ihrem Inhalt zu bewerten und Metriken nur ergänzend zu
nutzen.}
\begin{description}
\item[{\rot\bf Zwei Kennzahlen}]~\\[-\lblskip]
\bi
  \ite \textbf{Impact Factor} (Zeitschrift): Zitate eines Jahres auf die
    Artikel der zwei Vorjahre, geteilt durch deren Zahl
    \parencite{garfield2006history}
  \ite \textbf{h-Index} (Person): $h$ Arbeiten mit je mindestens $h$
    Zitaten \parencite{hirsch2005index}
  \ite beide zählen Zitate, also \emph{Aufmerksamkeit}
\ei
\os{\newpage}

\item[{\rot\bf Was Zitate nicht messen}]~\\[-\lblskip]
\bi
  \ite wenige Artikel tragen den Impact Factor einer Zeitschrift
    \parencite{seglen1997impact}
  \ite Bekannte werden häufiger zitiert: Matthäus-Effekt
    \parencite{merton1968matthew}
  \ite zitiert wird auch, was widerlegt oder zurückgezogen wurde
  \ite Fachkulturen zitieren unterschiedlich viel
  \ite DORA und Leiden Manifesto: Inhalt vor Kennzahl
    \parencite{dora2012,hicks2015leiden}
\ei
\os{\newpage}
\end{description}

\begin{merkbox}[Merke: Zitate]
Eine Zitation sagt: \enquote{Diese Arbeit wurde zur Kenntnis genommen.} Ob
sie zustimmend, kritisch oder als abschreckendes Beispiel zitiert wurde,
steht nicht in der Zahl.
\end{merkbox}
\os{\newpage}

%==============================================================================
\section{Laterales Lesen und SIFT}
%==============================================================================
\nsl{Wie prüfen Profis eine unbekannte Quelle? \textcite{wineburg2019lateral}
haben Historiker, Studierende einer Elite-Universität und professionelle
Faktenprüfer dieselben Websites bewerten lassen. Die Faktenprüfer waren
schneller und lagen häufiger richtig, obwohl sie weniger auf den Seiten
selbst lasen. Sie verließen die Seite fast sofort und schlugen in neuen
Tabs nach, wer hinter ihr steht. Wineburg und McGrew nennen das
\emph{laterales Lesen}, im Unterschied zum \emph{vertikalen} Lesen, bei dem
man die Seite von oben nach unten prüft und dabei den Gestaltungsmitteln
des Absenders ausgeliefert bleibt. Im Begleitbuch heißt es dazu: Die Website
ist die Bühne des Absenders, und wer auf der Bühne nach der Wahrheit sucht,
sucht dort, wo die Ablenkung am stärksten ist (\buch{Kap.~15b}).

Mike Caulfield hat das Vorgehen der Faktenprüfer zu vier Schritten
verdichtet, die er SIFT nennt \parencite{caulfield2023verified}. Sie
eignen sich für Zeitungsartikel und Websites ebenso wie als erste Stufe der
Prüfung wissenschaftlicher Quellen.}
\begin{description}
\item[{\rot\bf Vertikal oder lateral?}]~\\[-\lblskip]
\begin{center}
\resizebox{0.85\textwidth}{!}{%
\begin{tikzpicture}[font=\sffamily,>=Stealth,
  pg/.style={draw=rwuviolet,thick,fill=rwuviolet!8,minimum width=22mm,minimum height=30mm},
  tb/.style={draw=rwucyan,thick,fill=rwucyan40!35,minimum width=27mm,minimum height=12mm,align=center,font=\sffamily\small}]
  \node[pg] (v) at (0,0) {};
  \node[above] at (v.north) {\bfseries vertikal};
  \foreach \y in {1.0,0.5,0,-0.5,-1.0}{\draw[rwugray] (-0.8,\y) -- (0.8,\y);}
  \draw[->,very thick,rwuviolet] (1.4,1.2) -- (1.4,-1.2);
  \node[below,align=center,font=\sffamily\small] at (0,-1.7) {die Seite gründlich\\lesen};
  \node[pg] (l) at (6,0) {};
  \node[above] at (l.north) {\bfseries lateral};
  \foreach \y in {1.0,0.5}{\draw[rwugray] (5.2,\y) -- (6.8,\y);}
  \node[tb] (t1) at (10,1.7) {Wer steht\\dahinter?};
  \node[tb] (t2) at (10,0) {Was sagen\\andere?};
  \node[tb] (t3) at (10,-1.7) {Woher stammt\\die Zahl?};
  \draw[->,thick] (l.east) -- (t1.west);
  \draw[->,thick] (l.east) -- (t2.west);
  \draw[->,thick] (l.east) -- (t3.west);
  \node[below,align=center,font=\sffamily\small] at (6,-1.7) {die Seite früh\\verlassen};
\end{tikzpicture}}
\end{center}
\bi
  \ite Faktenprüfer lesen weniger und urteilen besser
    \parencite{wineburg2019lateral}
\ei
\os{\newpage}

\item[{\rot\bf SIFT in vier Schritten}]~\\[-\lblskip]
\begin{center}
\small
\begin{tabular}{@{}p{0.3\textwidth}p{0.6\textwidth}@{}}
\toprule
\textbf{S}top & innehalten: Kenne ich die Quelle? Was will ich wissen? \\
\textbf{I}nvestigate the source & außerhalb der Quelle nachlesen, wer sie ist \\
\textbf{F}ind better coverage & bessere oder unabhängige Berichte zur selben Behauptung suchen \\
\textbf{T}race claims & Behauptung, Zitat, Bild zum Ursprung zurückverfolgen \\
\bottomrule
\end{tabular}
\end{center}
\bi
  \ite nach \textcite{caulfield2023verified}
\ei
\os{\newpage}
\end{description}

\nsl{\textbf{Stop} heißt zweierlei: den Impuls bremsen, eine Information
sofort zu glauben oder weiterzugeben, und sich klarmachen, wozu man sie
braucht. Für eine Hausarbeit genügt nicht, was für ein Gespräch in der
Mensa genügt. \textbf{Investigate the source} heißt nicht, die
\enquote{Über uns}-Seite der Quelle zu lesen, sondern in einem neuen Tab
nachzuschlagen, was \emph{andere} über sie sagen; Wikipedia ist dafür ein
brauchbarer erster Halt. \textbf{Find better coverage} heißt: Wenn die
Quelle unklar bleibt, eine bessere zur selben Behauptung suchen, statt die
schlechte weiter zu analysieren. \textbf{Trace claims} heißt, eine Zahl,
ein Zitat oder ein Bild bis zum ursprünglichen Kontext zu verfolgen. Bei
Wissenschaftsmeldungen endet diese Spur im Idealfall bei der Studie selbst.
Dort zeigt sich oft, dass die Studie etwas Vorsichtigeres behauptet als die
Schlagzeile.}

%==============================================================================
\section{Prebunking und Debunking}
%==============================================================================
\nsl{Quellenkritik ist eine Fähigkeit. Sie wirkt aber nur, wenn man sie
anwendet, bevor sich eine Falschinformation festgesetzt hat. Die
Desinformationsforschung zeigt, dass widerlegte Behauptungen weiter auf
das Denken abfärben, selbst wenn die Korrektur verstanden und akzeptiert
wurde \parencite{lewandowsky2012misinformation}. Das Begleitbuch deutet
diesen \emph{Continued Influence Effect} mit dem Begriffspaar aus LE~1: Die
Korrektur erreicht den Belief, die Falschmeldung sitzt als Alief
(\buch{Kap.~15b}).

Daraus folgen zwei Strategien. Beim \emph{Debunking} wird eine bereits
verbreitete Falschbehauptung widerlegt; wirksam ist das vor allem, wenn die
Korrektur die Lücke mit einer Erklärung füllt, statt nur \enquote{falsch}
zu sagen \parencite{lewandowsky2020debunking}. Beim \emph{Prebunking}
werden Menschen vorab mit einer abgeschwächten Form der
Manipulationstechnik konfrontiert. Die Idee geht auf McGuires
Inokulationstheorie zurück \parencite{mcguire1964inoculation} und wurde in
Experimenten mit kurzen Videos, darunter einem Feldexperiment auf YouTube,
bestätigt \parencite{roozenbeek2022inoculation}. Schon die bloße Frage
\enquote{Stimmt das eigentlich?} verbessert, was Menschen teilen
\parencite{pennycook2021accuracy}. Der Staunmoment dieser Stunde war eine
solche Impfung: Sie haben die Wirkung der schönen Fälschung an sich selbst
erlebt.}
\begin{description}
\item[{\rot\bf Korrigieren reicht nicht}]~\\[-\lblskip]
\bi
  \ite widerlegte Behauptungen wirken weiter
    \parencite{lewandowsky2012misinformation}
  \ite die Korrektur erreicht den Belief, die Falschmeldung sitzt als Alief
    (\buch{Kap.~15b})
  \ite \textbf{Debunking}: die Lücke mit einer Erklärung füllen, nicht nur
    \enquote{falsch} sagen \parencite{lewandowsky2020debunking}
\ei
\os{\newpage}

\item[{\rot\bf Impfen statt heilen}]~\\[-\lblskip]
\bi
  \ite \textbf{Prebunking}: die Technik in schwacher Dosis vorab erleben
    \parencite{mcguire1964inoculation}
  \ite wirkt gegen \emph{Techniken}, nicht nur gegen einzelne Behauptungen
    \parencite{roozenbeek2022inoculation}
  \ite schon die Frage \enquote{Stimmt das?} hilft
    \parencite{pennycook2021accuracy}
  \ite der Zaubertrick ist Prebunking: Täuschung erleben, dann auflösen
\ei
\os{\newpage}

\item[{\rot\bf Skepsis ohne Zynismus}]~\\[-\lblskip]
\bi
  \ite Ziel ist \emph{kalibriertes} Vertrauen, nicht Misstrauen gegen alles
  \ite wer alles für inszeniert hält, landet bei Expertenverachtung
    \parencite{mede2020populism}
  \ite dass Betrug aufgedeckt wird, spricht \emph{für} die Wissenschaft
  \ite offen bleiben und prüfen \parencite{sagan1995demon}
\ei
\nsl{Diese Folie ist wichtig, weil Quellenkritik eine Nebenwirkung hat. Wer
gelernt hat, dass alles gefälscht sein \emph{kann}, schließt leicht, dass
alles gefälscht \emph{ist}. Das Begleitbuch setzt dagegen eine
Dosierungsregel: hohes Vertrauen dort, wo Herkunft, Methode und
unabhängige Bestätigung es rechtfertigen, und Zurückhaltung dort, wo sie
fehlen (\buch{Kap.~15b}). Carl Sagan hat diese Balance als Verbindung von
Offenheit und Skepsis beschrieben \parencite{sagan1995demon}.}
\os{\newpage}
\end{description}

%==============================================================================
\section{Übung in der Stunde}
%==============================================================================
\begin{description}
\item[{\rot\bf SIFT in Kleingruppen (30 Min.)}]~\\[-\lblskip]
\bi
  \ite drei Quellen: eine Studie, ein Zeitungsartikel, ein
    Wikipedia-Eintrag (Links auf dem Lernportal)
  \ite füllen Sie das SIFT-Raster aus (Übungsblatt, Aufgabe~2)
  \ite verfolgen Sie mindestens eine Zahl bis zur Originalstudie
  \ite Ergebnis: ein Gewicht je Quelle (hoch / mittel / gering) mit
    Begründung
\ei
\os{\newpage}

\item[{\rot\bf Quellenkritik in fünf Schritten}]~\\[-\lblskip]
\bi
  \ite \textbf{1 Innehalten}: Was will ich wissen, was löst die Quelle in
    mir aus?
  \ite \textbf{2 Herkunft und Absicht}: wer, mit welchem Interesse?
  \ite \textbf{3 Art und Methode}: Primärstudie, Review, Bericht? Wie
    gewonnen?
  \ite \textbf{4 Spur zum Original}: Zahlen und Zitate zurückverfolgen
  \ite \textbf{5 Gegenprobe}: unabhängige Bestätigung, Widerruf,
    Gegenposition
\ei
\os{\newpage}
\end{description}

\begin{lehrhinweis}
Die fünf Schritte entstehen in der Sicherung gemeinsam an der Tafel; die
Folie dient als Abgleich. Erwartungshorizont für die SIFT-Übung (abhängig
von den gewählten Quellen): \emph{Studie} -- Herkunft über Zeitschrift
und Institution prüfbar, Methode im Text offengelegt, Gewicht hoch bis
mittel je nach Design und Replikation; die Gruppe soll nachsehen, ob es
Korrekturen oder einen Widerruf gibt. \emph{Zeitungsartikel} -- Herkunft
über das Medium bestimmbar, aber Tertiärquelle; in Schritt~T zeigt sich
meist die Zuspitzung gegenüber der Studie. Gewicht: als Hinweis brauchbar,
als Beleg in einer Hausarbeit nicht. \emph{Wikipedia} -- Autorenschaft
anonym, aber Versionsgeschichte und Diskussionsseite einsehbar;
Einzelnachweise sind der Weg zur Primärquelle. Gewicht: guter
Ausgangspunkt für laterales Lesen, kein zitierfähiger Beleg für
Fachaussagen. Häufiger Fehler: Gruppen bewerten vertikal (\enquote{die
Seite wirkt seriös}). Dann nachfragen: \enquote{Was sagen \emph{andere}
über diese Quelle?}
\end{lehrhinweis}

%==============================================================================
\section{Lernziele und Aufgaben}
%==============================================================================
\begin{description}
\item[{\rot\bf Nach dieser Einheit können Sie}]~\\[-\lblskip]
\bi
  \ite Apophänie und Bestätigungsfehler erklären und an Beispielen zeigen
  \ite Quellen nach Herkunft, Intention, Methode, Peer Review und
    Aktualität bewerten
  \ite Primärstudie, Review, Meta-Analyse und graue Literatur
    unterscheiden
  \ite die Grenzen von Plausibilitätsurteilen und Zitationsmetriken
    benennen
  \ite eine Quelle lateral nach dem SIFT-Schema prüfen
\ei
\os{\newpage}

\item[{\rot\bf Bis zur nächsten Woche}]~\\[-\lblskip]
\bi
  \ite \textbf{SIFT-Raster} für die drei Quellen fertigstellen
    (Übungsblatt, Aufgabe~2)
  \ite \textbf{Diskussion} (max.\ 1 Seite): \enquote{Eine Studie hat 500
    Zitierungen. Was sagt das über ihre Qualität?} Strukturierte
    Gegenargumentation
  \ite \textbf{Tagebuch des Staunens}: ein Muster, das Sie diese Woche
    gesehen haben, und wie Sie geprüft haben, ob es echt ist
  \ite die eigene annotierte Bibliographie aus LE~3 mit den fünf Schritten
    nachprüfen
\ei
\os{\newpage}
\end{description}

\handoutinput{handouts/LE04}
